\subsection{The degree-normalized Bessel lowering scalar}
\label{archimedean-normalized-lowering}

We now retain the original normalization by subgroup periods.  This
normalization is an input to the original family; it is distinct from
the still separate quotient $L^2$ isometry.  Under this input the
normalized lowering scalar can be evaluated, and the preceding product
estimate gives its required uniform bound
\cite{normalized-lowering-source}.

Let $n=2r+k+1$, $k\ge1$, and let $H,H_s$ be the original unitary
representations with parameters $\mu,\mu-e_s$.  Retain their signed
dual-standard tensor decomposition $i_s,p_s$, with signs $\epsilon_s$,
as in Theorem~\ref{original-tensor-projection-bound}.  All operators
below act on common invariant smooth domains.  For each representation
write $\mathcal D$ for the linear span of group translates of its
compact finite smooth vectors.  This is the original period test
domain; it includes the compact finite vectors used in the norm
recursion.  Write $S$ for the
original Bessel subgroup, and $\mathcal B:H\to W$,
$\mathcal B_s:H_s\to W$ for its Bessel maps.  Assume $\mathcal B$
is nonzero on $\mathcal D$.

The normalization is the following exact assertion.  With the
original subgroup measure and character, the absolutely convergent
operator period satisfies
\begin{equation}\label{lowering-period-normalization}
 \mathcal P_H(w,v)=R(\mathcal B w,\mathcal B v),\qquad
 R(a,b)\xi=a\langle b,\xi\rangle,
\end{equation}
for smooth $w$ and $v\in\mathcal D$, and the same assertion holds
for every $H_s$.  The kernel defining $\mathcal P_H$ is the original
compact representation and unipotent character.  In particular all
periods in the following comparison use the same subgroup measures.
This is exactly the source's normalized Bessel-map condition, not a
choice to make a weighted integral have mass one.

Let $f_0$ be the original first flag vector, and let
$\omega=\operatorname{flat}(f_0)$ in the dual standard representation.
It is fixed by $S$.  Define
\[
 A_s v=p_s(\omega\otimes v),\qquad
 C_s y=\operatorname{ev}_{\Gamma}(\omega,i_sy),
\]
where $\Gamma$ is the original invariant Hermitian form on the dual
standard module.  The signed isometry gives the positive-Hilbert-space
adjoint identity $C_s=\epsilon_s A_s^*$.  Suppose the original lowering
scalar $\kappa_s$ is characterized by
\begin{equation}\label{lowering-scalar-definition}
                  \mathcal B_sA_s=\kappa_s\mathcal B
\end{equation}
on $\mathcal D$.  The tensor maps preserve these test domains:
apply an intertwiner to a translated compact finite tensor and use
finite dimensionality of the standard factor.  The existence of the Pieri decomposition
and of the Bessel family is retained; no different family is introduced.

\begin{lemma}\label{lowering-normalized-adjoint}
Under \eqref{lowering-period-normalization} and
\eqref{lowering-scalar-definition},
\[
                \mathcal B C_s=\epsilon_s\overline{\kappa_s}
                                                    \mathcal B_s.
\]
\end{lemma}
\begin{proof}
Because $\omega$ is $S$-fixed, $C_s$ intertwines the $S$ actions.
For every $g\in S$ its matrix coefficients therefore satisfy
\[
        \langle v,\sigma(g)C_sy\rangle
                   =\epsilon_s\langle A_sv,\sigma_s(g)y\rangle.
\]
Multiply by the unchanged period kernel and integrate.  The absolute
convergence in \eqref{lowering-period-normalization} justifies the
comparison and gives
\[
 R(\mathcal B C_sy,\mathcal B v)
       =\epsilon_s R(\mathcal B_s y,\mathcal B_sA_sv)
       =\epsilon_s\overline{\kappa_s}
                           R(\mathcal B_s y,\mathcal B v).
\]
Choose $v\in\mathcal D$ with $\mathcal B v\ne0$.  Applying both
sides to $\mathcal B v$ cancels the positive number
$\|\mathcal B v\|^2$ and proves the identity.  It also covers
$\mathcal B_s=0$, when \eqref{lowering-scalar-definition} forces
$\kappa_s=0$.
\end{proof}

We give the operator evaluation needed for the scalar, with the shifts
visible.  For ordered row and column lists of length $d$ use the
original Capelli minor
\[
 Q_{I,J}(u)=\sum_{\sigma\in\mathfrak S_d}\operatorname{sgn}(\sigma)
       \prod_{a=0}^{d-1}
       \bigl(E_{i_{\sigma(a)},j_a}+(u-a)\delta_{i_{\sigma(a)},j_a}\bigr),
\]
in the displayed order.  The empty minor is the identity.  The usual
matrix-unit commutator relation implies row and column alternation and
the ordered cofactor identity
\begin{equation}\label{lowering-quantum-comatrix}
 \mathcal C(Y)\bigl((Y-(n-1)/2)I-T\bigr)
             =Q_n(Y+(n-1)/2)I,
\end{equation}
where $(Tx)_a=-\sum_bE_{ab}x_b$ and
$\mathcal C(Y)_{ij}=(-1)^{i+j}Q_{\widehat j,\widehat i}(Y+(n-1)/2)$.
The determinant on the right acts on each coordinate.  To obtain
\eqref{lowering-quantum-comatrix}, append a row to a cofactor and use
the ordered end-row expansion; a repeated row contributes zero, and
the missing row contributes the full minor.  The appended row acts
on the right at parameter $Y-(n-1)/2$.  Thus no adjacent noncommuting
operators have been exchanged.  This is the same cofactor calculation
as for the original ordered minors
\cite{normalized-lowering-prior-minors}.

The original central character and compact Capelli parameter give
\begin{align}
 Q_n(Y+(n-1)/2)&=\prod_t(Y+\mu_t)I,                  \label{lowering-full-character}\\
 \mathcal B Q_{W,W}(v)
   &=\prod_{i=1}^k(v-k/2+1/2+u_i)\mathcal B,        \label{lowering-compact-character}
\end{align}
where $W=(r,r+1,\ldots,r+k-1)$ is the positive compact block in an
adapted Bessel frame.  These retain the original representation
parameters; the compact identity is part of the original side datum.

Let $P_s=i_sp_s$.  The tensor Casimir computation gives
$TP_s=(-\mu_s-(n-1)/2)P_s$.  On this summand,
\eqref{lowering-quantum-comatrix} and
\eqref{lowering-full-character} imply the polynomial identity
\[
              \mathcal C(Y)P_s=\prod_{t\ne s}(Y+\mu_t)P_s.
\]
One first cancels $Y+\mu_s$ away from its single zero, and then uses
equality of polynomial coefficients.  Consequently
\begin{equation}\label{lowering-projector-cofactor}
              \mathcal C(-\mu_s)=D_sP_s,
        \qquad D_s=\prod_{t\ne s}(\mu_t-\mu_s)\ne0.
\end{equation}

To evaluate the relevant entry, choose the original Witt frame with
the negative vector $z$ of norm $-T_0$, $T_0>0$, represented by
$\tau e_{r+k}$, where $\tau=\sqrt{T_0}$.  The outer isotropic
chains retain their original unit scale.  For an upper triangular
matrix unit outside the central block, the differentiated Bessel
character is a scalar $p_{ab}$, with
\[
 |p_{ab}|=
 \begin{cases}
 1,&b=a+1,\quad b<r\text{ or }a>r+k,\\
 \tau,&(a,b)=(r-1,r+k)\text{ or }(r+k,r+k+1),\\
 0,&\text{otherwise}.
 \end{cases}
\]
The two entries involving nonexistent indices are omitted when $r=0$.
This frame is obtained by choosing an orthonormal basis of the positive
orthogonal complement of $z$ in the original central Hermitian space;
it changes neither $T_0$ nor the flag vectors or the period measures.

In this frame the inclusion--projection matrix $P_s$ gives exactly
\[
 C_sA_sv=
 \begin{cases}
       (P_s(e_{n-1}\otimes v))_0,&r>0,\\
       -T_0(P_s(e_{n-1}\otimes v))_{n-1},&r=0.
 \end{cases}
\]
Indeed the original dual flag vector has coordinates $e_{n-1}$
and the evaluation selects coordinate $0$ when $r>0$.  For $r=0$
these coordinates are $-\tau e_{n-1}$ and $\tau$ times evaluation
at $n-1$.  This identifies the cofactor entry before taking its norm.

For $r>0$, the required cofactor has rows $0,\ldots,n-2$ and columns
$1,\ldots,n-1$.  Expand in its first row, putting that row factor on
the left.  Only its character entry survives after applying
$\mathcal B$.  Successively remove rows $0,\ldots,r-2$ with columns
$1,\ldots,r-1$, then row $r-1$ with column $r+k$.  The product of
these scalars has modulus $\tau$, and the minor parameter is unchanged.
The remaining rows are $r,\ldots,n-2$; its columns are
$r,\ldots,r+k-1,r+k+1,\ldots,n-1$.

Next put the last remaining column first and expand in that column,
again putting its factor on the left.  Only its last-row character
entry survives.  Remove columns $n-1,n-2,\ldots,r+k+1$ in this way.
Column alternation contributes only signs of modulus one.  Each
first-column expansion lowers the parameter by one.  These $r$
steps contribute a second factor of modulus $\tau$ and leave
\[
                    Q_{W,W}(u-r).
\]
The scalar terms on the removed diagonal positions are zero: at each
row step the removed row is absent from the column list, and at each
column step the removed column is absent from the row list.  Thus
these are valid left eigenfunctional substitutions, not substitutions
inside a noncommuting product.  At $u=-\mu_s+(n-1)/2$ the surviving
parameter is exactly $-\mu_s+k/2$.

For $r=0$ no row or column is removed.  The cofactor is already the
compact minor at $-\mu_s+k/2$.  Here $f_0=z$; the dual tensor vector
and its evaluation contribute respectively $-\tau$ and $\tau$.
Their product has modulus $T_0$, the same factor as the two strings
when $r>0$.

Combining this calculation with
\eqref{lowering-compact-character} and
\eqref{lowering-projector-cofactor} gives a scalar $c_s$ with
\begin{equation}\label{lowering-evaluation-scalar}
 \mathcal B C_sA_s=c_s\mathcal B,
 \qquad
 |c_s|=T_0\frac{\prod_i|\mu_s-1/2-u_i|}
                     {\prod_{t\ne s}|\mu_s-\mu_t|}.
\end{equation}
Lemma~\ref{lowering-normalized-adjoint} and
\eqref{lowering-scalar-definition} also give
$c_s=\epsilon_s|\kappa_s|^2$.  Nonvanishing of $\mathcal B$ justifies
this scalar equality; no quotient-norm nonvanishing is used.

\begin{theorem}\label{original-normalized-lowering-bound}
Retain the preceding original period normalization, tensor family,
characters and frame, and the parameter hypotheses of
Theorem~\ref{original-lowering-product-bound}.  If
$b\ge B\ge10^6n^7$, then
\[
            |\kappa_s|^2\frac{d_\mu}{d_{\mu-e_s}}\le T_0b^4,
\]
where both formal degrees use the same original Haar measure.
\end{theorem}
\begin{proof}
For a fixed group and Haar measure the formal-degree constant is
independent of the parameter, so the previously established
Harish--Chandra product formula gives
\[
 \frac{d_\mu}{d_{\mu-e_s}}
       =\frac{\prod_{t\ne s}|\mu_s-\mu_t|}
                    {\prod_{t\ne s}|\mu_s-1-\mu_t|}.
\]
All denominators are nonzero by the original gap hypothesis.  Multiply
this identity by \eqref{lowering-evaluation-scalar}, use
$|c_s|=|\kappa_s|^2$, and apply
Theorem~\ref{original-lowering-product-bound}.
\end{proof}

The theorem reconstructs the original normalized lowering-scalar
supplier from the original subgroup-period normalization.  It does
not prove existence of that normalized Bessel family, its identification
with the original Jacquet transform, or the quotient $L^2$ mass bound.
Those analytic obligations remain necessary for the leading-minor
norm theorem.
